% Chapter opener (rewritten 2026-10-31, Aida-style single paragraph; refactored 2026-11-06 to nine
% stages, one per section, in \input order).
% Order of \input (no main file yet):
%   01a_inhouse_tracer      tracer on all BL scans (cited tool, not own)     sec:method-tracer
%   01b_medis_reference     Medis collapsed to full arteries; tracer vs Medis sec:method-medis
%   02_population_outcome   population, outcome                              sec:method-population
%   02b_segment_split       dividing the traced arteries into segments       sec:method-split
%   03_prediction_model     pooled cloglog GLMM, base, LR test, CV           sec:method-models, sec:method-cv
%   04_model_size           parameter budget and feature selection          sec:method-size, sec:method-selection
%                           (Riley criteria: technical_background/03b)
%   05_features             umbrella, then 05a_curvature, 05b_bifurcation_angle,
%                           05c_daughter_ratio, 05d_taper                     sec:method-features
%   06_feature_validity     measurement properties                          sec:method-validity
%   07_curvature_fpca       secondary analysis: curvature shape (FPCA)       sec:method-fpca
% Deleted 2026-11-06: 10_medis_validation.tex (superseded by 01b_medis_reference.tex) and
%   11_cross_validation.tex (moved into 03_prediction_model.tex).
% Retired, not \input: retired/04_feature_selection.tex (merged into 04_model_size.tex).
% Overview figure still to make (style guide: an overview figure opens the methods).
\providecommand{\clarify}[1]{\textcolor{red}{\textbf{Clarification needed:} #1}}
\chapter{Methodology}
\label{ch:methodology}

The methodology followed in this work is summarized in Figure~\ref{fig:method-overview} and consists of nine main stages:
(1) extraction of the coronary centerlines and lumen contours from all baseline scans with an in-house tracer;
(2) comparison of the extracted centerlines and lumen with Medis as a reference, as no ground truth was available;
(3) selection of the plaque-free baseline cohort and definition of the outcome from the visual analysis;
(4) division of the traced arteries into the segments of the visual analysis;
(5) the prediction model and its cross-validation;
(6) the number of parameters the cohort can support and the choice of geometric features within it;
(7) definition of the geometric features;
(8) assessment of their measurement properties, against synthetic shapes and known groups, and of the variance explained by known risk factors;
and (9) a secondary, descriptive analysis of the shape of the curvature profile.

\clarify{overview figure of the nine stages (label fig:method-overview).}
